\section{Erasure-Coded History}
\label{sec:history}
\subsection{State versus history}
The pivotal design decision is to treat \emph{state} and \emph{history} differently.
\emph{State}, account balances, nonces, contract storage, is the materialized result of
replaying history; it is bounded (one entry per active account or storage slot) and every
validator keeps it in full locally, exactly as Bitcoin keeps the UTXO set and Ethereum the
state trie. \emph{History}, the block bodies that produced that state, is unbounded and is
the \emph{only} thing fragmented away. Because state is derivable, losing a node's local copy
is not data loss: it is recomputed by replay, or fetched as a snapshot from a peer and checked
against the committed chain. This separation is precisely what lets ``no node stores all
history'' coexist with fast, deterministic validation: validation reads \emph{state}, which is
local; only audit, sync, and dispute resolution need \emph{history}, which is recoverable.

\subsection{Encoding, placement, reconstruction}
On commit, a block body $b$ of size $|b|$ is RS-encoded into $T=K{+}M$ shards, each of size
$\lceil |b|/K\rceil$ (data shards padded; parity shards equal size). The erasure shape
$(K,M)$ is chosen adaptively from the live node count $n$ (\S\ref{sec:overview}); the leader
stamps it into the signed header, so peers read the shape rather than agreeing on it. Each
shard index $s$ is assigned to the top-$r$ nodes under $\mathrm{HRW}(\textit{blockID},s)$,
where $r$ is a replication factor that also scales with $n$. A node persists only the shards
it wins and discards $b$. The header commits $\mathit{BodyHash}(b)$ and the vector of per-shard
hashes; any node gathering $K$ valid shards reconstructs $b$ and verifies
$\mathit{BodyHash}=\,$commitment, so a lying holder is detected and skipped.

\emph{Concrete walkthrough.} Consider a $128$-node network producing block $1000$ with a $1$\,MiB
body. The shape is K16\,M8, so the body becomes $24$ shards of $64$\,KiB each, and the header
commits one $32$-byte hash per shard plus the body hash. Node $X$ computes
$\mathrm{HRW}(\textit{id}_{1000},s)$ for every shard $s$ and finds it ranks in the top
$r{=}3$ for, say, shards $\{3,7\}$; it stores those two ($128$\,KiB) and discards the
$1$\,MiB body. Months later an auditor requests block $1000$: it queries holders, collects any
$16$ of the $24$ shards (each verified against its committed hash), runs RS decode to rebuild the
$1$\,MiB body, and checks the body hash. If eight holders are offline the remaining sixteen still
suffice; if a malicious holder returns a corrupted shard, its hash mismatch excludes it and another
holder's shard is used instead. Node $X$ never reconstructed the block to validate the chain ahead
of it; it only did so on this explicit audit request.

\begin{figure}[t]
\centering
\begin{tikzpicture}[font=\scriptsize, node distance=4.5mm,
  lcell/.style={draw, rounded corners, fill=gray!8, align=center, minimum width=20mm, minimum height=6mm},
  a/.style={-{Stealth[length=1.6mm]}}]
  \node[lcell] (s1) {1. leader\\assembles txs};
  \node[lcell, below=of s1] (s2) {2. apply state\\transition (local)};
  \node[lcell, below=of s2] (s3) {3. sign header\\(BodyHash, shardHashes)};
  \node[lcell, below=of s3] (s4) {4. broadcast block};
  \node[lcell, right=14mm of s4] (s5) {5. peers apply\\state transition};
  \node[lcell, above=of s5] (s6) {6. RS-encode body};
  \node[lcell, above=of s6] (s7) {7. HRW: keep mine,\\drop full body};
  \node[lcell, above=of s7] (s8) {8. on demand:\\fetch $K$, reconstruct};
  \draw[a] (s1)--(s2); \draw[a] (s2)--(s3); \draw[a] (s3)--(s4);
  \draw[a] (s4)--(s5); \draw[a] (s5)--(s6); \draw[a] (s6)--(s7); \draw[a] (s7)--(s8);
\end{tikzpicture}
\caption{Block lifecycle: state is applied and kept by all; the body is encoded, only assigned
shards are retained, and old bodies are reconstructed on demand.}
\label{fig:lifecycle}
\end{figure}

\subsection{Storage model}
\label{sec:storagemodel}
Let each shard have $r$ holders and let the $T$ shards be spread over $n$ nodes by HRW. The
expected number of shard copies a node holds for one block is $Tr/n$, each of size $|b|/K$, so
expected per-node storage per block is
\[
S_{\text{node}} \;=\; \frac{Tr}{n}\cdot\frac{|b|}{K}.
\]
Full replication stores $|b|$ per node, giving a reduction factor
\[
\rho \;=\; \frac{|b|}{S_{\text{node}}} \;=\; \frac{Kn}{Tr}.
\]
With ShadowLedger's policy ($T,K$ saturate at $K{=}16,M{=}8$ and $r{=}3$ for large $n$),
$\rho \approx Kn/(Tr) = 16n/(24\cdot3)=n/4.5$, i.e.\ per-node history cost falls as
$O(1/n)$ while full replication is constant. Below $n\!\approx\!4$ the fixed parity overhead
$T/K$ dominates and fragmentation is not yet a win, an honest small-network regime we report
rather than hide. \S\ref{sec:eval} validates this model empirically against rendezvous
placement.

\subsection{Sync and state snapshots}
\label{sec:sync}
A node that joins late, or that lost its local state, must reach the current head without
necessarily replaying every historical body. Two paths exist. The \emph{header-first} path
downloads the (small, signed) header chain and verifies its links and signatures, giving the node
a trustworthy skeleton of heights, commitments, and the validator set; bodies are then fetched and
reconstructed only as needed (\S\ref{sec:eval} quantifies that cost). The \emph{snapshot} path
fetches a materialized state at a recent height from a peer and adopts it as a \emph{checkpoint}:
the node sets its replay base (\S\ref{sec:fork}) to that height, so it can still reorg above the
checkpoint but trusts the snapshot below it. Because the snapshot is the fold of all prior blocks,
its correctness is in principle checkable by replay; making it succinctly \emph{verifiable}
without replay needs a committed state root, which the present header omits (\S\ref{sec:disc}). In
both paths, steady-state operation then reverts to keeping only HRW-assigned shards: sync cost is
paid once, storage savings persist.
